\section{Clearance strips transverse to physical test curves}
\label{sec:cf-transverse}
A measure bound alone gives no short-curve control. The positive
separation between an unvisited disk and the incident disks supplies
a transverse derivative which remains nonzero even as an incident
angle approaches grazing. We specify the time orientation explicitly:
forward near misses are tested on increasing-slope curves, and
\emph{past-flight} near misses on decreasing-slope curves.

Use physical arclength $s=R\alpha$ increasing counterclockwise and
$p=\sin\varphi$. Fix a finite $U>0$ and put
\begin{equation}\label{eq:cf-curve-constants}
 \begin{split}
 \kappa_-&=1/R_+,\qquad \kappa_+=1/R_-,\\
 D_0&=\kappa_++U+(1-R_+)^{-1},\\
 \ell_0&=\pi/(2D_0),\qquad
 a_0=\kappa_-(g_*-\delta_*),\\
 C_{\rm geo}&=2N_*\sqrt{1+U^2}/a_0.
 \end{split}
\end{equation}
All lengths below are Euclidean arclength in $(s,\varphi)$, not
configuration-space arclength. A test is a connected graph over an
$s$ interval, of length at most $\ell_0$. Let $E_\delta^+$ be the
forward physical near-miss event $\{N_\delta\ge1\}$, and let
$E_\delta^-$ be the analogous event on the preceding actual flight.

\begin{theorem}[A grazing-uniform, one-flight strip bound]
\label{thm:cf-transverse}
For every $C^1$ test graph $W^+$ with
$0\le d\varphi/ds\le U$, and every $C^1$ test graph $W^-$ with
$-U\le d\varphi/ds\le0$,
\begin{equation}\label{eq:cf-strip}
 |W^+\cap E_\delta^+|\le C_{\rm geo}\delta,
 \qquad |W^-\cap E_\delta^-|\le C_{\rm geo}\delta.
\end{equation}
The bound is uniform in $R$, in the homogeneity strip of the test,
and in its distance to the grazing boundary. Arbitrary Borel
restrictions, including exact record restrictions, can be added to
the sets on the left. Regular pieces are understood with their
physical multiplicity one.
\end{theorem}
\begin{proof}
On $W^+$ the direction angle is $\vartheta=s/R+\varphi(s)$.
For a fixed candidate center $z\ne0$ set
\[
 t_z=(z-q)\cdot v_\vartheta,\qquad
 b_z=(z-q)\cdot v_\vartheta^\perp.
\]
Differentiation in the increasing $s$ coordinate gives the signed
identity
\begin{equation}\label{eq:cf-transverse-derivative}
 b'_z=-\cos\varphi-t_z\vartheta',\qquad
 \kappa_-\le\vartheta'\le\kappa_++U.
\end{equation}
At a physical near miss, Lemma~\ref{lem:cf-foot} yields
$t_z\ge g_*-\delta_*$. Hence $-b'_z\ge a_0$ there. It is the
$t_z\vartheta'$ term, not a lower bound on $\cos\varphi$, which
prevents degeneration.

We must not multiply this estimate by an uncontrolled number of
components. Choose a continuous argument $\beta(s)=\arg(z-q(s))$
on the graph. Since $|z-q|\ge1-R_+$, one has
$|\beta'|\le(1-R_+)^{-1}$. Moreover
\[
 (\beta-\vartheta)'\le (1-R_+)^{-1}-\kappa_-<0,
\]
because $R_+<1/2$. Its total variation is at most
$D_0|W^+|\le\pi/2$. Thus the set $t_z>0$, equivalently
$\cos(\beta-\vartheta)>0$, is at most one interval: two distinct
positive-cosine intervals are separated by a negative interval of
angular length $\pi$. On the interval $t_z>0$, formula
\eqref{eq:cf-transverse-derivative} says that $b_z$ strictly decreases.
Each of the two impact bands $(R,R+\delta)$ and
$(-R-\delta,-R)$ therefore has multiplicity at most one in this
coordinate. Restricting to the portions which are actual near misses
can break them into pieces, but cannot increase that multiplicity.
By the one-dimensional change of variables on this monotone map,
\[
 \int_{\{s:\ z\text{ is a physical near miss}\}}ds
                 \le 2\delta/a_0.
\]
Since $d\ell\le\sqrt{1+U^2}\,ds$, summing over the finite candidate
set from Lemma~\ref{lem:cf-foot} proves the first bound. This sum
counts only centers in one physical flight. It does not count
iterated collision words or homogeneous continuity components.

For the second bound, the collision time-reversal involution is
$\mathcal I(s,\varphi)=(s,-\varphi)$. The outgoing velocity at
$\mathcal I x$ reverses the incoming velocity at $x$. Consequently
its next straight flight is the reversal of the previous actual
flight at $x$. This proves
\begin{equation}\label{eq:cf-reversal}
 E_\delta^-=\mathcal I E_\delta^+=T_RE_\delta^+
\end{equation}
on the regular collision space. The involution is an isometry in
$(s,\varphi)$ and sends $W^-$ to an increasing-slope graph.
Apply the first bound. Borel restrictions decrease lengths. Taking
monotone limits of truncated regular test pieces gives the same
bound at grazing limits; an exactly tangent parameter contributes
zero transverse length by the strict derivative above. This last
statement does not delete any paired-source boundary trace in the
older coarea construction.
\end{proof}

\begin{corollary}[Bounded multipliers on stable tests]
\label{cor:cf-multiplier}
If $|h|\le H$ pointwise on regular collision states, then for every
$W^-$ as above and every bounded Borel test $\psi$,
\[
 \left|\int_{W^-}\one_{E_\delta^-}h\psi\,d\ell\right|
       \le H\|\psi\|_\infty
                   \min\{|W^-|,C_{\rm geo}\delta\}.
\]
This is a bound for bounded functions with specified representatives.
It is not an assertion that multiplication by the indicator is
bounded on an entire space of anisotropic distributions.
\end{corollary}
\begin{proof}
Bound the integrand by $H\|\psi\|_\infty\one_{E_\delta^-}$ and
use Theorem~\ref{thm:cf-transverse} together with the trivial length
bound. Representatives on singular sets are fixed to zero.
\end{proof}
